% texte.tex — f166r, vv. 295–336 (conversion initiale des travaux antérieurs)
\vers[n=294]{}{}{δούλοις γὰρ οὐκ ἔξεστ᾽ ἀκούειν τῶν λόγων.}{car il n'est pas permis aux esclaves d'ouïr ce qui se dit ici.}

\vers[n=295, indent=0.5em, findent=0.3em, metre={prose}]{ΚΟ.}{Coryphée}{εὐφημία ἔστω, |}{\didasc{La prêtresse :} Recueillez-vous, recueillez-vous.}
\vers[indent=0.5em, findent=0.3em]{}{}{εὐφημία ἔστω. |}{Priez les Thesmophores,}
\vers[n=296, indent=0.5em, findent=0.3em]{}{}{εὔχεσθε ταῖν Θεσμοφόροιν |}{et Ploutos et Calligénie,}
\vers[n=297, indent=0.5em, findent=0.3em]{}{}{καὶ τῷ Πλούτῳ καὶ τῇ Καλλιγενείᾳ |}{et la Nourricière,}
\vers[indent=0.5em, findent=0.3em]{}{}{καὶ τῇ Κουροτρόφῳ |}{et Hermès et les Charites,}
\vers[n=300, indent=0.5em, findent=0.3em]{}{}{καὶ τῷ Ἑρμῇ καὶ ταῖς Χάρισιν |}{de rendre cette ecclésie et présente réunion}
\vers[n=301, indent=0.5em, findent=0.3em]{}{}{ἐκκλησίαν τήνδε καὶ σύνοδον τὴν νῦν |}{la meilleure et la plus avantageuse qu'il se pourra,}
\vers[n=302, indent=0.5em, findent=0.3em]{}{}{κάλλιστα καὶ ἄριστα ποῆσαι, |}{pour le plus grand profit de la cité des Athéniens}
\vers[indent=0.5em, findent=0.3em]{}{}{πολυωφελῶς μὲν τῇ πόλει τῇ Ἀθηναίων, |}{et pour notre propre bien.}
\vers[n=305, indent=0.5em, findent=0.3em]{}{}{τυχηρῶς δ᾽ ἡμῖν αὐταῖς. |}{Que celle qui aura agi et parlé}
\vers[n=306, indent=0.5em, findent=0.3em]{}{}{καὶ τὴν δρῶσαν καὶ ἀγορεύουσαν |}{le mieux en faveur du peuple des Athéniens}
\vers[n=307, indent=0.5em, findent=0.3em]{}{}{τὰ βέλτιστα περὶ τὸν δῆμον τὸν Ἀθηναίων |}{et du peuple des femmes voie prévaloir ses avis.}
\vers[indent=0.5em, findent=0.3em]{}{}{καὶ τὸν τῶν γυναικῶν |}{Demandez aux dieux cette faveur, et pour vous-mêmes toutes les prospérités.}
\vers[suite]{ΚΟ.}{Coryphée}{ταύτην νικᾶν. |}{Demandez aux dieux cette faveur, et pour vous-mêmes toutes les prospérités.}
\vers[n=310, indent=0.5em, findent=0.3em]{}{}{ταῦτ᾽ εὔχεσθε καὶ ὑμῖν αὐταῖς τἀγαθά. |}{Io Péan ! Io Péan ! Soyons en Joie.}
\vers[n=311, indent=0.5em, findent=0.3em]{}{}{ἰὴ παιὼν ἰὴ παιὼν ἰὴ παιών. χαίρωμεν.}{}
\vers[n=312, indent=2em, findent=1.2em, metre={lyr.}]{ΧΟ.}{Chœur}{δεχόμεθα καὶ θεῶν γένος}{Nous acceptons ! Et nous supplions}
\vers[n=313, indent=2em, findent=1.2em]{}{}{λιτόμεθα ταῖσδ᾽ ἐπ᾽ εὐχαῖς}{la gent des dieux d'apparaître}
\vers[n=314, indent=2em, findent=1.2em]{}{}{φανέντας ἐπιχαρῆναι.}{et de prendre en gré ces prières.}
\vers[n=315, indent=2em, findent=1.2em]{}{}{Ζεῦ μεγαλώνυμε χρυσολύρα τε}{Zeus au grand renom, et toi, dieu à la lyre d'or,}
\vers[n=316, indent=2em, findent=1.2em]{}{}{Δῆλον ὃς ἔχεις ἱεράν,}{qui tiens Délos la sainte,}
\vers[n=317, indent=2em, findent=1.2em]{}{}{καὶ σὺ παγκρατὲς κόρα}{et toi, vierge toute-puissante,}
\vers[n=318, indent=2em, findent=1.2em]{}{}{γλαυκῶπι χρυσόλογχε πόλιν}{aux yeux glauques, à la lance d'or, qui habites}
\vers[n=319, indent=2em, findent=1.2em]{}{}{οἰκοῦσα περιμάχητον, ἐλθὲ δεῦρο.}{une cité enviée, viens ici. Et toi aussi,}
\vers[n=320, indent=2em, findent=1.2em]{}{}{καὶ πολυώνυμε θηροφόνη}{vierge aux noms multiples, meurtrière de bêtes,}
\vers[n=321, indent=2em, findent=1.2em]{}{}{Λατοῦς χρυσώπιδος ἔρνος,}{rejeton de Léto aux yeux d'or.}
\vers[n=322, indent=2em, findent=1.2em]{}{}{σύ τε πόντιε σεμνὲ Πόσειδον}{Et toi, auguste Poseidon, thalassien,}
\vers[n=323, indent=2em, findent=1.2em]{}{}{ἁλιμέδον}{modérateur de la mer,}
\vers[n=324, indent=2em, findent=1.2em]{}{}{προλιπὼν μυχὸν ἰχθυόεντα,}{quitte le gouffre poissonneux}
\vers[suite]{ΧΟ.}{Chœur}{οἰστροδόνητον}{fouetté des vents.}
\vers[n=325, indent=2em, findent=1.2em]{}{}{Νηρέος εἰναλίου τε κόραι}{Et vous, filles de Nérée le marin,}
\vers[n=326, indent=2em, findent=1.2em]{}{}{Νύμφαι τ᾽ ὀρείπλαγκτοι.}{et vous, Nymphes des montagnes.}
\vers[n=327, indent=2em, findent=1.2em]{}{}{χρυσέα τε φόρμιγξ}{Que la phorminx d'or}
\vers[n=328, indent=2em, findent=1.2em]{}{}{ἰαχήσειεν ἐπ᾽ εὐχαῖς}{accompagne nos prières, et puissions-nous}
\vers[n=329, indent=2em, findent=1.2em]{}{}{ἡμετέραις· τελέως δ᾽ ἐκκλησιάσαιμεν Ἀθηνῶν}{mener à bonne fin notre ecclésie,}
\vers[n=330, indent=2em, findent=1.2em, metre={lyr.}]{}{}{εὐγενεῖς γυναῖκες.}{femmes de pure lignée athénienne.}
\vers[n=331]{ΚΟ.}{Coryphée}{εὔχεσθε τοῖς θεοῖσι τοῖς Ὀλυμπίοις}{Invoquez les dieux et les déesses de l'Olympe,}
\vers[n=332]{}{}{καὶ ταῖς Ὀλυμπίαισι καὶ τοῖς Πυθίοις}{les dieux et les déesses de Pytho,}
\vers[n=333]{}{}{καὶ ταῖσι Πυθίαισι καὶ τοῖς Δηλίοις}{les dieux et les déesses de Délos}
\vers[n=334]{}{}{καὶ ταῖσι Δηλίαισι τοῖς τ᾽ ἄλλοις θεοῖς,}{et les autres dieux.}
\vers[n=335]{}{}{εἴ τις ἐπιβουλεύει τι τῷ δήμῳ κακὸν}{Si quelqu'un forme quelque dessein contre le peuple}
\vers[n=336]{}{}{τῷ τῶν γυναικῶν ἢ 'πικηρυκεύεται}{des femmes, fait des ouvertures}
